\problemname{Bartering}

We have all heard the fascinating story of how someone managed to trade a paperclip for a pen, then traded the
pen for a door handle, and so on, until they eventually traded their way to a house. 
In today's housing market, you think this sounds like a brilliant idea and decide to try something similar.

After analyzing the market, you realize that there are $N$ distinct items and $M$ possible trades. 
A trade consists of a set of items that you can exchange for exactly one other item. 
For example, a trade could be 2 apples and 3 pears for an orange. Formally, a trade can be described as 
$$
C_{i,0} \cdot F_0 + C_{i,1} \cdot F_1 + \dots + C_{i,N-1} \cdot F_{N-1} \implies 1 \cdot F_{r_i}.
$$
Here, $C_{i,j}$ is the coefficient for $F_j$ in trade $i$. This means that $C_{i,j}$ is the number of $F_j$
required for the trade. $r_i$ is the item you receive as a result of trade $i$. 
Note that not all items need to be part of a particular trade; for these items, $C_{i,j} = 0$.

Each trade can be performed any number of times. Your goal is to calculate how many paperclips are needed to
acquire a house.

The item $F_0$ represents the paperclip, and the item $F_{N-1}$ represents the house.

\section*{Input}
The first line contains two integers $N$ and $M$ ($2 \leq N \leq 10^5$, $1 \leq M \leq 10^5$), the number of
items and the number of trades.

Next, the $M$ trades are described. Each description begins with two integers $r_i$ ($0 \leq r_i < N$)
and $k_i$ ($1 \leq k_i \leq N$),
where $r_i$ is the item you receive from trade $i$, and $k_i$ is the number of distinct items used in the trade.

Then follow $k_i$ lines, each containing two integers $C_{i,j}$ ($1 \leq C_{i,j} \leq 100$) and $j$ ($0 \leq j < N$),
which indicate that the coefficient $C_{i,j}$ belongs to item $F_j$ in the trade. Note that 
items with coefficient $0$ are not included in the input.

Let $K$ be the sum of all $k_i$ values across all trades. It holds that $K \leq 10^5$.

It is guaranteed that all items can be obtained by starting with a paperclip, and you can purchase each item for
at most $10^9$ paperclips each.

\section*{Output}
Print a single integer: the minimum number of \emph{paperclips} required to trade your way to a \emph{house}.

\section*{Scoring}
Your solution will be tested on a set of test groups, each worth a number of points. Each test group contains
a set of test cases. To get the points for a test group you need to solve all test cases in the test group.

\noindent
\begin{tabular}{| l | l | p{12cm} |}
  \hline
  \textbf{Group} & \textbf{Points} & \textbf{Constraints} \\ \hline
  $1$ & $12$ & $N = 3$ \\ \hline
  $2$ & $10$ & $k_i = 1$ \\ \hline
  $3$ & $25$ & There is an order of items such that for every item
  $f$, all trades where $f$ is received only contain items that come earlier in the order than $f$. \\ \hline
  $4$ & $26$ & $K \leq 1000$ \\ \hline
  $5$ & $27$ & No additional constraints. \\ \hline
\end{tabular}
